%% ══════════════════════════════════════════════════════════
%%  ITEM — Shuffled Time   Sponsorēšanas piedāvājums  (Latviski)
\newcommand{\sponsor}{BALCIA}
%% ══════════════════════════════════════════════════════════

\documentclass[a4paper,9.5pt]{article}

\usepackage[top=0cm,bottom=0cm,left=0cm,right=0cm]{geometry}
\usepackage{xcolor}
\usepackage[T1]{fontenc}
\usepackage[utf8]{inputenc}
\usepackage{tikz}
\usepackage{graphicx}
\usepackage{tabularx}
\usepackage{parskip}
\usepackage{microtype}
\usepackage{lmodern}
\usepackage{textcomp}
\usepackage{enumitem}
\usepackage{multicol}
\usepackage{ifthen}
\usepackage{hyperref}

\hypersetup{hidelinks}
\pagestyle{empty}
\setlength{\parskip}{1.5pt}
\setlength{\columnsep}{0.8em}

\definecolor{raacream}{HTML}{F5F5DC}
\definecolor{raared}{HTML}{CC2222}
\definecolor{raadark}{HTML}{0D0D0D}
\definecolor{itemgold}{HTML}{D4A84B}
\definecolor{mygray}{HTML}{555555}

\newcommand{\sectionrule}[1]{%
  \vspace{3pt}%
  {\color{raared}\bfseries\normalsize #1}\\[-3pt]%
  {\color{raared}\rule{\linewidth}{0.8pt}}%
  \vspace{2pt}%
}

%% Performances karte: {datums}{nosaukums}{mākslinieks}{attēls/none}{apraksts}{bio}
\newcommand{\perfcard}[6]{%
  \noindent\begin{minipage}[t]{0.37\linewidth}
    \vspace{0pt}
    {\small\color{raared}\textbf{#1}}\\[1pt]
    {\bfseries\small #2}\\[2pt]
    {\footnotesize\textit{#3}}\\[5pt]
    \ifthenelse{\equal{#4}{none}}{%
      {\setlength{\fboxsep}{0pt}%
       \colorbox{black!8}{\parbox[c][3.6cm][c]{\linewidth}{%
         \centering\footnotesize\color{mygray}attēls drīz}}}%
    }{%
      \makebox[\linewidth][c]{%
        \includegraphics[height=3.6cm,width=\linewidth,keepaspectratio=true]{#4}}%
    }%
  \end{minipage}%
  \hfill%
  \begin{minipage}[t]{0.60\linewidth}
    \vspace{0pt}
    {\small #5}%
    \par\vspace{5pt}%
    {\footnotesize\color{mygray}\textbf{Biogrāfija:}\enspace #6}%
  \end{minipage}%
  \par\vspace{0.35cm}%
  {\color{raared!35}\rule{\linewidth}{0.4pt}}%
  \par\vspace{0.35cm}%
}

\begin{document}

%% ══════════════════════════════════════════════════════════
%%  1. LAPPUSE — PILNEKRĀNA PLAKĀTS + PIEDĀVĀJUMA SATURS
%% ══════════════════════════════════════════════════════════
\noindent
\begin{tikzpicture}[x=1cm,y=1cm]
  \useasboundingbox (0,0) rectangle (21,29.7);

  %% Plakāts — izstiepts līdz lapas malām, bez atkāpēm
  \node[anchor=south west, inner sep=0pt] at (0,0) {%
    \includegraphics[width=21cm,height=29.7cm,keepaspectratio=false]{img/item-poster.jpg}%
  };

  %% Tumšs pārklājums teksta salasāmībai
  \fill[black,opacity=0.67] (0,0) rectangle (21,29.7);

  %% ── GALVENE ──────────────────────────────────────────────

  \node[anchor=north west, inner sep=0pt] at (0.9,29.4) {%
    
\includegraphics[height=1.1cm]{img/raa_logo_v2.png}%
  };

  \node[anchor=north east, inner sep=0pt] at (20.1,29.5) {%
    \fontsize{5.5}{7}\selectfont\color{raacream!50}%
    \textls[180]{S P O N S O R Ē Š A N A S \ \ P I E D Ā V Ā J U M S}%
  };

  \node[anchor=center] at (10.5,26.8) {%
    \fontsize{26}{31}\selectfont\bfseries\color{raacream}%
    ITEM\,—\,SHUFFLED TIME%
  };

  \node[anchor=center] at (10.5,25.9) {%
    \fontsize{9}{11}\selectfont\color{raacream!80}%
    \textit{5.–11. jūlijs 2026\enspace·\enspace RAA.SPACE, Matīsa iela 8, Rīga}%
  };

  \node[anchor=center] at (10.5,24.85) {%
    \fontsize{9}{11}\selectfont\color{itemgold}%
    Iesniegts\enspace\textbf{\sponsor}%
  };

  \fill[raared] (0,24.15) rectangle (21,24.45);

  %% ── SATURS ───────────────────────────────────────────────
  \node[anchor=north west, text width=19cm, inner sep=0pt] at (1.0,24.0) {%
    \begin{minipage}[t]{19cm}
      \color{raacream}
      \setlength{\parskip}{1.5pt}
      \setlength{\columnsep}{0.8em}

      \vspace{0.18cm}
      \sectionrule{Par pasākumu}

      \small
      ITEM\,—\,Shuffled Time ir septiņu dienu performatīva nedēļa RAA.SPACE Rīgā.
      Katru dienu mākslinieks, māksliniece vai mākslinieku grupa izmantos RAA galveno
      zāli (138\,m\textsuperscript{2} plaša telpa ar desmit logiem uz Matīsa ielu), lai
      24 stundu laikā iedzīvinātu, veidotu, piepildītu un atklātu telpu no jauna.
      Performance var ilgt visas 24 stundas, vai arī īsāku laika posmu, ietverot vienu
      vai vairākas uzstāšanās.
      Programma apvieno septiņas mākslinieciskas balsis ilgstošās performances,
      ķermeniskās mākslas, eksperimentālās dzejas, aromagrāfijas un interaktīvās
      instalācijas žanros. Organizē RAA, atbalsta Veldze Radošā Rūpnīca.

      \vspace{0.1cm}

      \begin{multicols}{2}

        \sectionrule{Programma}

        {\footnotesize\setlength{\parskip}{0pt}
        \begin{tabular}{@{}p{1.35cm}p{5.65cm}@{}}
        \textbf{Sv. 5.07}   & \textbf{Ilze Aulmane}\\
                             & \textit{(sīkāka informācija drīz)}\\[1pt]
        \textbf{P. 6.07}    & \textbf{NOW} — Maruta Gredzena\\
                             & Interaktīva velšanas performance\\[1pt]
        \textbf{O. 7.07}    & \textbf{N.3} — Ernests Ā\\
                             & Virve \& izturība, 18:00–24:00\\[1pt]
        \textbf{T. 8.07}    & \textbf{Rijība} — K.\,Šuksta \& S.\,Murins\\
                             & Aromagrāfijas performance\\[1pt]
        \textbf{C. 9.07}    & \textbf{Presence in Absence Pt.\,3}\\
                             & Lisette Ros\\[1pt]
        \textbf{Pk. 10.07}  & \textbf{Time Scores} — Vanessa Brzič\\
                             & Fluxus participatorā performance\\[1pt]
        \textbf{S. 11.07}   & \textbf{Laiks Telpā} — Ilze Mazpane\\
                             & \& Jana Zariņa\\
                             & Ilgstoša mūzikas-ķermeņa-telpas performance\\
        \end{tabular}
        }

        \columnbreak

        \sectionrule{Auditorija un sasniedzamība}

        {\small
        \begin{itemize}[leftmargin=*,itemsep=1pt,topsep=0pt]
          \item \textbf{200+} apmeklētāji nedēļas laikā
          \item Starptautiski un vietējie mākslinieki, mākslas profesionāļi,
                ziņkārīgā publika
          \item Instagram \textbf{@raa\_riga} \& \textbf{@veldze.radosarupnica}:\\
                kopā \textbf{7\,000+} sekotāji
          \item Mediju sasniedzamība: \textbf{100gkultura},
                \textbf{riga-neighbourhood},\\
                \textbf{arterritory.com}, \textbf{notboring.riga}
        \end{itemize}
        }

        \vspace{0.05cm}
        \sectionrule{Kāpēc tas ir svarīgi}

        {\small Ilgstoša un dzīvā performance ir neapstrādāta un neatkārtojama.
        \textbf{\sponsor} atbalstītu ITEM — nedēļu, kas notiek RAA.SPACE, telpā ar
        stabilu performances un vizuālās mākslas programmu kopš tās dibināšanas:
        Homo Novus, Kefīrs Studio, 2 Annas, Alīna Burlakova, Andrejs Osokins un
        daudzi citi. Sponsorēšana šeit ir patiess kultūras ieguldījums — redzams un
        jūtams iesaistītai, mākslai atvērtai auditorijai.}

      \end{multicols}

      \vspace{0.05cm}

      \sectionrule{Sponsorēšanas pakete — \texteuro\,2\,000}

      {\small
      \begin{tabularx}{19cm}{@{}>{\bfseries}p{3.3cm}X@{}}
      Pieminēts visur        & \sponsor{} tiek pieminēts katrā pasākumā, visos drukas
                               materiālos un RAA.SPACE pasākuma lapā.\\[2pt]
      Sociālie mediji        & Veltītas publikācijas un stāsti @raa\_riga,
                               @veldze.radosarupnica, riga-neighbourhood, notboring.riga,
                               100gkultura, arterritory.com (7\,000+ kopējā sasniedzamība).\\[2pt]
      VIP pieeja             & 3 rezervētas vietas katru vakaru. Ielūgumi uz mākslinieku
                               tikšanās reizēm un atklāšanas vakariem.\\[2pt]
      Uzstāšanās iespēja     & Īss uzstāšanās vārds atklāšanas vakarā (5. jūl.).\\[2pt]
      Klātbūtne uz vietas    & \sponsor{} ekspozīcija RAA foajē visas nedēļas garumā.\\[2pt]
      Preses atspoguļojums   & Nosaukts kā galvenais sponsors visos mediju paziņojumos.\\
      \end{tabularx}
      }

    \end{minipage}%
  };

  %% ── APAKŠDAĻAS JOSLA ─────────────────────────────────────
  \fill[raadark,opacity=0.88] (0,0) rectangle (21,0.82);
  \node[anchor=center] at (10.5,0.41) {%
    \fontsize{7.5}{9}\selectfont\color{raacream!80}%
    \textbf{RAA.SPACE}\enspace·\enspace
    Matīsa iela 8, Rīga\enspace·\enspace
    \href{https://raa.space}{raa.space}\enspace·\enspace
    \href{mailto:enter@raa.space}{enter@raa.space}\enspace·\enspace
    @raa\_riga%
  };

\end{tikzpicture}

%% ══════════════════════════════════════════════════════════
%%  2.+ LAPPUSE — PERFORMANCES KARTES
%% ══════════════════════════════════════════════════════════
\newpage
\newgeometry{top=1.5cm,bottom=1.5cm,left=1.5cm,right=1.5cm}

\noindent
{\color{raared}\bfseries\large ITEM\,—\,SHUFFLED TIME}%
{\small\color{mygray}\quad Performances programma · 5.–11. jūlijs 2026 · RAA.SPACE, Rīga}%
\\[-4pt]%
{\color{raared}\rule{\linewidth}{0.8pt}}

\vspace{0.45cm}

%% ── 1 · Ilze Aulmane ─────────────────────────────────────
\perfcard%
  {Svētdiena, 5. jūlijs}%
  {Ilze Aulmane}%
  {Ilze Aulmane}%
  {img/item-aulmane.jpg}%
  {Pilns nosaukums, apraksts un vizuālais materiāls tiks publicēts drīz.}%
  {Biogrāfija drīz.}%

%% ── 2 · NOW ──────────────────────────────────────────────
\perfcard%
  {Pirmdiena, 6. jūlijs}%
  {NOW}%
  {Maruta Gredzena}%
  {img/item-now.jpg}%
  {Interaktīvajā performancē \textit{Now} Maruta Gredzena pēta, vai komunikācija
   bez vārdiem var būt patiesa. Apmeklētāji kļūst par dalībniekiem, izvēloties un
   uzvelkot vieglus, papīra bāzētus velšanas skulptūras savā pašreizējā noskaņojumā
   — mīkstas, plūksnainas, ģeometriskas vai akmeņveida formas veido pagaidu otro
   ādu. Katra izvēle pārveido mainīgu kolektīvo kompozīciju, kur katrs kļūst gan
   par skatītāju, gan par mākslas darbu. Vai godīgums var pastāvēt bez valodas?
   Papīra skulptūru instalācija ar objektiem un strēmelēm no griestiem veido
   telpas fonu.}%
  {Maruta Gredzena interesējas par to, kas notiek, kad nekas nepaliek mierā.
   Dzīve pastāvīgo pārcelšanos dēļ mainījusi viņas praksi no objektu radīšanas
   uz saskarsmes telpu veidošanu. Saknes porcelāna un keramikas skulptūrā,
   paplašinot caur grafikas mākslu, viņas darbs atklājas caur instalācijām, kas
   izmanto papīru, materiālu un rotaļīgumu, lai pētītu komunikāciju ārpus valodas.
   IG:~@em.gaa}%

%% ── 3 · N.3 ──────────────────────────────────────────────
\perfcard%
  {Otrdiena, 7. jūlijs}%
  {N.3}%
  {Ernests Ā}%
  {img/item-n3.jpg}%
  {\textit{Saistošā pārliecība par divu pretstatu esamību vienmēr liek apjukumā
   maldīties.}\par\vspace{3pt}%
   Ernests ieradīsies ar virvi, ko vispirms aptīs ap vienu no galerijas kolonnām.
   Tās galu piesies sev pie kājas un ies apkārt kolonnai, cenšoties attālināties,
   metot lokus gan telpā, gan ārpus tās sienām — līdz spēku izsīkumam vai ja
   viņu apturēs ārējie faktori. Ja izdosies atbrīvoties, iespējams, aptinās virvi
   no jauna. Īsi pārtraukumi performanci nepārtrauc. Performances laikā Ernests
   nerunās. Ieplānotais laiks: 18:00–24:00.}%
  {Savā praksē Ernests Ā nododas rotaļīgai realitātes ikdienišķo un nedefinējamo
   aspektu izpētei neierobežotās izteiksmes formās.}%

%% ── 4 · Rijība ───────────────────────────────────────────
\perfcard%
  {Trešdiena, 8. jūlijs}%
  {Rijība}%
  {Kristiāna Šuksta \& Sandris Murins}%
  {img/item-gluttony.jpg}%
  {Eksperimentāla aromagrāfijas performance, kas savieno divpadsmit vārdus —
   katru ar savu īpašu aromātu — nepārtraukti mainot secību un tempu. Darbs atklāj,
   kā smarža var stiprināt, pretoties vai pārveidot teksta nozīmi, pētot vēlmi,
   atmiņu un gaļas patēriņa sensorās dimensijas.}%
  {Sandris Murins ir latviešu mākslinieks, Smaržas kameras izgudrotājs un Smaržas +
   Mākslas festivāla dibinātājs — popularizē aromagrāfiju kā topošu mākslas valodu.
   Kristiāna Šuksta ir latviešu dzejniece, performances māksliniece un festivāla
   Starptelpa tehniskā producente. Viņas prakse apvieno dzeju un tiešraides performanci,
   pētot personisko pieredzi un cilvēciskās attiecības ar asprātību un humoru.}%

%% ── 5 · Presence in Absence, Piece Three ─────────────────
\perfcard%
  {Ceturtdiena, 9. jūlijs}%
  {Presence in Absence, Piece Three}%
  {Lisette Ros (holandiešu)}%
  {img/item-presence.jpg}%
  {\textit{Presence in Absence} ir Lisette Ros sērija, ko viņa uzsāka Mehiko
   2023. gadā, pētot \textit{Prombūtnē esošo ķermeni} — klātbūtnes uztveri
   prombūtnē. Trešajā daļā Ros saskaras ar savu klātbūtni, uzņemot dabisku
   hipnotisku vielu, kas kopš 2014. gada ir daļa no viņas naktsvakara rituāla
   saistībā ar DSPS (aizkavētas miega fāzes sindroms). Darbs manipulē ar kontroli
   un apziņu reāllaikā — skatītāju priekšā.}%
  {Lisette Ros (dz. 1991, Nīderlande) ir holandiešu konceptuālā un performances
   māksliniece ar indonēziešu saknēm. Viņas darbs pēta identitāti kā sociālu
   konstruktu caur ķermeni kā pētījuma rīku un mediju. Starptautiski prezentēts:
   Venēcijas Arhitektūras biennāle, SXSW, Van Abbemuseum, Gētes institūts,
   MWoods Pekina, Casa Lü CDMX. 2021. gadā rezidence pie Marinas Abramović.}%

%% ── 6 · Time Scores ──────────────────────────────────────
\perfcard%
  {Piektdiena, 10. jūlijs}%
  {Time Scores}%
  {Vanessa Brzič (horvātu)}%
  {img/item-timescores.jpg}%
  {Apmeklētāji izvelk kartes ar īsām instrukcijām — piemēram ,,šķēlē'',
   ,,sasienu'', ,,apgāz'', ,,klausies'' — un tās veic telpā ar dažādiem
   priekšmetiem. Katra darbība maina telpu nākamajai. Tā veidojas dzīva, mainīga
   vide, kur katra atkārtotā karte iegūst jaunu nozīmi.}%
  {Vanessa Brzič ir horvātu starpdisciplināra māksliniece un OIST QISS pētniece,
   kas strādā mākslas un fundamentālo zinātņu krustpunktā. Sīkāka informācija drīz.}%

%% ── 7 · Laiks Telpā ─────────────────────────────────────
\perfcard%
  {Sestdiena, 11. jūlijs}%
  {Laiks Telpā}%
  {Ilze Mazpane \& Jana Zariņa}%
  {img/item-laikstelpa.jpg}%
  {Ilgstoša performance, kas esošajā RAA telpā, kāda tā ir, no uz vietas
   pieejamajiem objektiem radīs muzikāli-ķermeniski-telpisku subjektīvā laika
   pieredzi. Darbs pēta to, kā laiks izplešas, saraujas, skrien, apstājas un
   pazūd — bez pulksteņlaika kā vadotnes vai ierobežojuma.}%
  {Ilze Mazpane ir performances māksliniece no Rīgas ar izglītību un pieredzi
   filosofijā, vizuālajā mākslā, eksperimentālajā teātrī un radošajā rakstīšanā.
   Viņa veido no vietas izrietošas ilgstošās performances ar poētisku jutīgumu pret
   telpu un brīvību no pulksteņlaika tirānijas.
   Jana Zariņa ir vēsturisko taustiņinstrumentu un skaņas māksliniece, kuras darbība
   aptver seno mūziku, improvizāciju, elektroakustisko performanci un starpdisciplināro
   mākslu. Ievērojamākie projekti: elektroakustiskais duets \textit{maiak}, latviešu
   opera \textit{Ērces} (2024) un starptautiski baroka mūzikas pasākumi. Uzstājusies
   ar Latvijas Nacionālo simfonisko orķestri.}%

\end{document}
